\section{Structure-Exploiting Optimization}
\label{sec:algorithm}

\subsection{The declared $K=2$ special case}
If every event must contain exactly two rows, the world reduces to a
core-cover matching: every core has degree one, every buffer has degree at most
one, and buffer--buffer edges are excluded. Pair-additive queries become
weighted matching objectives. This is the Chicago benchmark and a useful
lower-complexity special case, but it cannot represent sequential
higher-cardinality events such as Figure~\ref{fig:chain}.

\begin{theorem}[$K=2$ core-cover matching equivalence]
\label{thm:k2matching}
Let $G=(C_0\cup B_0,E)$ contain every admissible two-row event incident to at
least one core row, with no buffer--buffer edge.  Feasible $K=2$ worlds are in
bijection with matchings in $G$ that saturate every vertex of $C_0$.  Under the
bijection, every pair-additive aggregate equals the corresponding edge-weight
sum.  Hence both frontier endpoints are exactly minimum- and maximum-weight
core-saturating matching problems.
\end{theorem}

\paragraph{Proof.}
Map each two-row event in a feasible world to its edge.  Disjoint event
membership implies that the selected edges form a matching, and core coverage
implies that every core vertex is saturated.  Conversely, every matching that
saturates $C_0$ defines disjoint two-row events covering each core exactly once;
unmatched buffers are unused.  These maps are inverse.  Pair additivity gives
equality of the world aggregate and selected-edge weight term by term.

The matching representation also localizes the source of a nondegenerate
frontier. For a matching $M$, write $w(M)=\sum_{e\in M}w_e$.

\begin{theorem}[Alternating decomposition of a fixed-cardinality $K=2$ frontier]
\label{thm:k2alternating}
Let $M^-$ and $M^+$ be any minimum- and maximum-weight feasible matchings of
the same cardinality $q>0$. Their symmetric difference is a vertex-disjoint
union $\mathcal A$ of alternating paths and even alternating cycles. Every
cycle contains equally many edges of $M^-$ and $M^+$; every path has edge-count
imbalance in $\{-1,0,1\}$, and these path imbalances sum to zero. Moreover,
\begin{equation}
 \overline H-\underline H
 =\frac{1}{q}\sum_{A\in\mathcal A}
 \left(\sum_{e\in A\cap M^+}w_e-
       \sum_{e\in A\cap M^-}w_e\right).
 \label{eq:alternating-width}
\end{equation}
Thus frontier width is generated by additive exchanges along alternating
components; equal cardinality may couple paths whose individual edge-count
imbalances are nonzero.
\end{theorem}

\paragraph{Proof.}
Every vertex has degree at most two in $M^-\mathbin{\triangle}M^+$ and is
incident to at most one edge of each matching. Each nontrivial connected
component is therefore a path or cycle whose edge colors alternate. A cycle is
even and balanced. An alternating path has equal color counts or one extra edge
of one color. Since $|M^-|=|M^+|=q$, the path imbalances sum to zero. Common
edges cancel from $w(M^+)-w(M^-)$; grouping the remaining terms by component
and dividing by $q$ proves Eq.~\eqref{eq:alternating-width}.

This theorem does not claim that every component can be exchanged alone. A
balanced cycle or path can be, but imbalanced paths must be exchanged in a
cardinality-balanced collection. Nor is the decomposition unique when an
endpoint has multiple optimal witnesses; endpoint values remain unique even
when the returned witness structure is not.

The equivalence also yields a support-omission sensitivity object that does not
depend on released geography.  Fix a declared base edge set $E_0\subseteq E$
and let $d_e=|e\cap C_0|$ for $e\in E\setminus E_0$, with $d_e=0$ otherwise.
For integer $\Gamma\in[0,|C_0|]$, impose
\begin{equation}
  \sum_{e\in E}d_e x_e\leq \Gamma.
  \label{eq:k2omission}
\end{equation}

\begin{proposition}[Nested candidate-omission frontier]
\label{prop:k2omission}
The feasible matching sets in Eq.~\eqref{eq:k2omission} are nested in
$\Gamma$.  At $\Gamma=0$ the set equals the core-saturating matchings using
only $E_0$; at $\Gamma=|C_0|$ it equals the unrestricted set on $E$.
Consequently every lower endpoint is weakly nonincreasing and every upper
endpoint weakly nondecreasing in $\Gamma$.
\end{proposition}

\paragraph{Proof.}
Nesting follows because only the right-hand side increases.  Every omitted
edge has positive core incidence, giving the first endpoint identity.  In any
core-saturating matching, summing $|e\cap C_0|$ over selected edges counts each
core exactly once and therefore equals $|C_0|$; the constraint is redundant at
the second endpoint.  Optimization over nested feasible sets gives the stated
endpoint monotonicity.

\begin{proposition}[When support contraction changes no decision]
\label{prop:k2thresholdcontraction}
Let $E'\subseteq E$ and suppose both fixed-$q$ feasible sets are nonempty. Then
\[
 \underline H_E\leq\underline H_{E'}\leq
 \overline H_{E'}\leq\overline H_E.
\]
For a threshold $\eta$ that is ambiguous under $E$, contraction to $E'$ leaves
the decision ambiguous exactly when
$\underline H_{E'}<\eta\leq\overline H_{E'}$. Hence strict contraction of a
continuous frontier need not improve decision certification.
\end{proposition}

\paragraph{Proof.}
Every feasible matching on $E'$ is feasible on $E$, so minimization over the
smaller set weakly raises the lower endpoint and maximization weakly lowers the
upper endpoint. The threshold rule certifies positive iff the lower endpoint
is at least $\eta$, certifies negative iff the upper endpoint is below $\eta$,
and is ambiguous otherwise, giving the equivalence.

\subsection{A compact connectivity characterization}
Fix exact intervals, a root row $r$, and a candidate span $[a,b)$. Let the
sorted distinct endpoints inside the span induce elementary open segments.
For a selected set containing $r$, positive-overlap connectivity with union
span $[a,b)$ is equivalent to:
\begin{enumerate}[leftmargin=*,itemsep=1pt,topsep=2pt]
\item every elementary segment is covered by at least one selected interval;
\item every internal endpoint boundary is strictly crossed by at least one
selected interval.
\end{enumerate}
The second condition excludes components that only touch at an endpoint.
Capacity adds an upper bound of $C$ to each segment-coverage row.

Interleave segment and boundary rows in temporal order. Each interval column
then contains ones in one consecutive block. This yields the pricing primitive.

\begin{theorem}[Integral rooted single-event oracle]
\label{thm:oracle}
For fixed exact intervals, root $r$, span $[a,b)$, capacity $C$, and additive
row weights, the linear relaxation of the segment-coverage, boundary-bridge,
capacity, root, and forced-companion formulation has an integral optimum.
Hence a minimum- or maximum-weight rooted event can be found by linear
programming. Enumerating observed endpoint spans and, only when needed, one
forced companion gives a polynomial-time exact single-event oracle.
\end{theorem}

\paragraph{Proof sketch.}
The interleaved segment/boundary incidence matrix has the consecutive-ones
property by columns. Signed duplicate rows for lower and upper coverage bounds,
plus identity rows fixing the root and companion, preserve total
unimodularity. The right-hand side is integral, so every extreme point is
integral. Candidate spans need only use observed interval endpoints. Appendix
details and exhaustive small-library comparisons are included in the artifact.

\subsection{Dantzig--Wolfe decomposition}
Let $\R_C$ be all feasible exact-time event columns. For event $R$, let
$b_R=|R\cap B_0|$ and introduce $\lambda_R\ge0$. The full relaxation for maximum
reachable buffer support is
\begin{align}
\max_{\lambda}\quad & \sum_{R\in\R_C}b_R\lambda_R
\label{eq:masterobj}\\
\text{s.t.}\quad
&\sum_{R\ni i}\lambda_R=1 && i\in C_0,\label{eq:coremaster}\\
&\sum_{R\ni j}\lambda_R\le1 && j\in B_0.\label{eq:buffermaster}
\end{align}
Core equalities and buffer inequalities couple otherwise independent event
columns. Given master duals, reduced cost is additive across selected rows, so
Theorem~\ref{thm:oracle} gives exact pricing. A phase-one artificial-core
master establishes feasibility.

\begin{proposition}[Full-master LP certificate]
\label{prop:cg}
If every rooted pricing call is solved exactly and the terminal minimum reduced
cost is nonnegative, the current restricted-master value equals the optimum of
Eqs.~\eqref{eq:masterobj}--\eqref{eq:buffermaster} over all feasible event
columns.
\end{proposition}

Local integrality does not transfer to the coupled master. The appendix gives a
capacity-two instance with
\[
\operatorname{OPT}_{LP}=4>3=\operatorname{OPT}_{IP}.
\]
A restricted integer solve over columns generated only for the root LP can
therefore miss the full integer optimum.

\subsection{Branch-compatible exact integer decomposition}
\label{sec:bp}
We close the support-maximization integer master with branch-and-price. At a
node, optional buffers may be fixed to total usage zero or one. Once row usage
is integral but event columns remain fractional, Ryan--Foster branching chooses
two active rows and creates:
\[
\text{together: both rows have identical membership in every event},\qquad
\text{separate: no event contains both.}
\]
For a priced root, together and separate restrictions are expanded into a
finite set of forced-in/forced-out cases. Each case adds only identity fixes
and variable bounds to the fixed-span formulation, so Theorem~\ref{thm:oracle}
still solves it exactly. The number of cases may grow exponentially with branch
depth; we make no polynomial-time claim for the full integer problem.

\begin{proposition}[Integer certificate]
\label{prop:bp}
Suppose each node relaxation is closed by exact branch-compatible pricing,
every child exhaustively implements its row-usage or Ryan--Foster disjunction,
and the branch queue terminates. Then the incumbent is a globally optimal
integer solution of the full event-column master. If computation stops early,
the best incumbent and largest open-node LP bound form a valid global lower and
upper bound.
\end{proposition}

\paragraph{Justification.}
The two children partition the integer solutions compatible with the parent.
Exact pricing makes each node bound valid for the full, not merely restricted,
column set. Standard branch-and-bound dominance then proves the claim. The
implementation replays every incumbent, records open-node bounds, and never
labels a time-limited cell optimal.

\subsection{Existential timestamp completion}
For set-valued timestamps, a world is feasible when there \emph{exists} one
exact interval choice inside each support that satisfies connectivity and
capacity. Replacing a support by its entire outer envelope is neither a valid
inner nor outer model for ordered events: it can add bridges while also
creating artificial concurrency. Our continuous mixed-integer model therefore
chooses latent starts and ends jointly with event membership. Each requested
support count receives one of four statuses: certified feasible witness,
proven infeasible, unresolved without incumbent, or unresolved invalid
incumbent. Only the first two are treated as resolved.
